% ---------------------------------------------------------------------------
\section{Introduction}

\subsection{The manuscript and the question}

The Voynich manuscript (Yale, Beinecke MS~408) is a codex of about 240 pages with drawings of unidentified plants,
astronomical and zodiacal diagrams, female figures in pools, pharmacy jars and a long final section of text arranged in
short paragraphs. Its vellum has been radiocarbon-dated to 1404--1438 at 95\% probability
\citep{uarizona2011,zyats2016}. The script has never been read. Prescott Currier showed in the 1970s that the pages fall
into two varieties of writing, now called languages A and B, and that more than one hand wrote them
\citep{currier1976}; Lisa Fagin Davis later distinguished five scribes \citep{davis2020}.

A century of proposed readings has not produced one that the scholarly community accepts. In recent years the debate
has narrowed to a question that statistics seems well placed to answer: is the text a meaningful message (a natural
language, perhaps enciphered) or a meaningless product of a procedure? Some studies find the text statistically close
to meaningless writing \citep{rugg2004,schinner2007,timmschinner2020,gaskellbowern2022}; others show that historically
plausible ciphers can produce Voynich-like text from Latin or Italian \citep{bowerngaskell2022,greshko2025,kinnison2026}.

\subsection{What statistics can decide}

Every statistical study of the manuscript measures regularities. This paper asks whether regularities, however many and
however strong, can decide the question at all. Steganography gives a precise answer in principle: if a text that
carries a message and a text that carries none are samples of the same distribution, no test on the text can tell them
apart \citep{cachin1998,hopper2002}. The question then becomes constructive. Can one build a book that reproduces the
measured writing rules of the Voynich manuscript and that carries, or does not carry, a message, with the two cases
indistinguishable by construction?

\subsection{Contributions}

\begin{enumerate}
\item \textbf{The voynichizer} (Sections~\ref{sec:generator} and~\ref{sec:evaluation}). A program that turns any text
  into a 207-folio Voynich-like book in EVA transliteration and recovers it exactly with the key
  \citep{voynichizer}. The text is compressed, encrypted and authenticated; its bits drive an arithmetic decoder that
  samples how many times each recurring Voynich word occurs on each page. Everything else (layout, invented words, word
  order) follows the manuscript's statistics and carries no information. To our knowledge, no earlier tool combines a
  cover model of a whole historical book, distribution-preserving embedding and authenticated encryption, evaluated
  against the measured properties of that book.
\item \textbf{A measured specification} (Sections~\ref{sec:rules}--\ref{sec:scorecard}). The scorecard and the judges
  that the generator must satisfy come from 716 logged experiments on how the manuscript is written. Each measure is
  compared with the same measure in 24 natural languages, constructed languages, volunteers' gibberish, published
  generators and ciphers and, systematically, 79 medieval manuscripts transcribed at the level of letter forms. Many of
  the regularities were observed before, and we cite them property by property (Section~\ref{sec:related}); our
  contribution there is to quantify them with the same code against the same comparison texts and to add a few
  properties that we have not found in the literature.
\item \textbf{A consequence} (Section~\ref{sec:discussion}). The generated book carries an encrypted message and comes
  close to the manuscript on the measures we know how to make. The regularities of the Voynich text therefore constrain
  how it was written, but they cannot by themselves show whether it says anything.
\end{enumerate}

\noindent The full research record is public: preregistrations, code, results, the laboratory notebook and a log of
all experiments with their outcomes \citep{repository}. Appendix~\ref{app:checks} lists the provisional readings that
later checks corrected, including those prompted by an external review of the first version of this paper.

% ---------------------------------------------------------------------------
\section{Related work}
\label{sec:related}

\subsection{Statistical regularities of the script}

\paragraph{Glyphs and words} The low conditional entropy of the next glyph was noted by \citet{bennett1976} and
quantified against hundreds of texts by \citet{lindemannbowern2020}. Words follow a rigid internal grammar
\citep{stolfi2000,zattera2022}; \citet{zattera2022} also found ``separable'' words that split into two attested words.
\citet{feaster2020} predicted mid-line spaces from glyph adjacencies with over 95\% accuracy and observed that uncertain
spaces fall where his rules are inconsistent; \citet{rozanova2026} describe two regimes of spaces, show that uncertain
spaces are physically narrower in the images and predict spaces from the flanking glyph pair with AUC 0.98.
\citet{matlach2022} revisit the positional roles of glyphs.

\paragraph{Adjacent words} \citet{currier1976} noted that in language B the end of a word influences the beginning of
the next, with \eva{qo-} words following words in \eva{-y}. Emma May Smith described the preference for \eva{o-} forms
after words in \eva{-n} and related transformations \citep{smith2017o,smith2017lastfirst}, and \citet{smithponzi2019}
measured the dependency between last and first glyphs across spaces (\eva{y.q} and \eva{n.o} frequent across spaces and
almost absent within words), concluding that word breaks are real boundaries. \citet{reddy2011} found that the last
letters of a word help predict the first letters of the next; \citet{caruana2022} and \citet{parisel2026a} studied
dependencies between adjacent words. \citet{feaster2022} warned that part of the \eva{y.q} and \eva{n.q} anomalies might
come from position in the line.

\paragraph{The line} \citet{currier1976} called the line a functional unit, noted line-final symbols and observed no
repeated sequence across a line break. \citet{vogt2012} measured word-length effects along the line;
\citet{smith2015linestart,smith2016linestart} described \eva{s-} and other forms added at line start, and
\citet{smith2016m} proposed \eva{-m} as a line-final form of \eva{-r}. \citet{feaster2021,feaster2022} defined
rightward and downward indices of glyphs (for example \eva{qo-} lies further left than \eva{o-}, \eva{Sh} further left
than \eva{ch}); \citet{bowernlindemann2021} summarise line and paragraph effects, and \citet{rozanova2026} find the
link between glyphs at independence across line breaks.

\paragraph{Lines above and below} \citet{timm2014} and \citet{timmschinner2020} observed that many words resemble words
written shortly before, often in the lines just above, and built a generator on this idea. On the left margin,
participants of the Voynich Ninja forum reported that a line-initial glyph tends to avoid the glyph directly above it
\citep{tavie2024}.

\paragraph{Paragraphs, pages, hands} Paragraph-initial gallows, especially \eva{p}, are well known
\citep{currier1976,dimperio1978,stafford2022,feaster2022}. \citet{davis2020} identified five hands, and
\citet{davis2025} noted that in the herbal section hands follow the bifolia, which she calls highly atypical; the
current binding is not the original order. Currier's languages have been confirmed with modern methods
\citep{farrugia2022,parisel2026b}. Immediate repetition of words has been measured by \citet{boxer2022} (about 1\% of
word pairs) and \citet{bowernlindemann2021}.

\subsection{Generators and ciphers}

Meaningless Voynich-like text has been produced with a Cardan grille \citep{rugg2004,zandbergen2021grille}, with the
self-citation procedure of \citet{timmschinner2020}, with Trithemius' word tables \citep{hermes2022} and with
recent generators that target statistical fingerprints: two amateur generators \citep{whitehatnetizen2026} and
\emph{voynich-fingerprint} \citep{sachak2026}, which reproduces 44 statistics on held-out folios. On the cipher side,
\citet{bowerngaskell2022} show that word-level metrics are compatible with some kinds of encipherment, and the Naibbe
cipher \citep{greshko2025} turns Latin or Italian into Voynich-like text by hand; its author presents it as a proof of
concept, not as the method used for the manuscript. \citet{kinnison2026} shows that Naibbe ciphertext reproduces the
manuscript's positional entropy profile.

\subsection{Steganography and its evaluation}

Transforming a text so that it imitates the statistics of another goes back to the mimic functions of
\citet{wayner1992}. \citet{cachin1998} defined the security of steganography as the distance between the distributions
of covers with and without a message, and \citet{hopper2002} gave provably secure constructions. Feeding encrypted bits
to an arithmetic decoder in order to sample from a language model, so that the message-carrying text follows the
model's distribution, is the method of \citet{ziegler2019}; \citet{kaptchuk2021} made it cryptographically secure.
Training a classifier to separate two samples and reading its accuracy is a classifier two-sample test
\citep{lopezpaz2017}. The voynichizer applies these ideas with a cover model of one specific historical book and is
evaluated against measured properties of that book.
